%======================================================================
\subsection{BW-PDM-06 --- Plan: the BOM per revision, where-used, and the push to Ops (24 Oct 2026, approved 24 Oct)}\label{sec:pdm06-plan}
%======================================================================

\eref{PDM-06} makes the vault know what every assembly revision is made of. Two things follow from that: the indented
BOM of a revision can be shown, exported and sent to Ops, and \emph{where-used} can answer which assembly revisions
contain a part. The acceptance stays as written: \emph{where-used of a bolt lists every assembly revision that uses it.}
It depends on the vault client (\eref{PDM-04}, building) and on the assembly BOM (\eref{ASM-04}, in progress). A spike
in this build (below) showed that the reference model can already explode a BOM and walk it backwards, but it keeps one
BOM per \emph{item}, not per revision, and nothing writes it yet.

\paragraph{Already decided.}
\begin{center}
\small
\begin{tabularx}{\linewidth}{@{}lX@{}}
\toprule
Decision (when) & What it means here\\
\midrule
Data model (19 Oct, figure~\ref{fig:pdm-model}) & A BOM line points from a parent \emph{revision} to a child \emph{item}, with
  quantity and find number. A released assembly picks up the child's latest released revision unless it is pinned.\\
One part per file (19 Oct, \eref{ASM-01}) & Every BOM line is a file with its own P/N; bodies inside a part are not
  BOM lines.\\
Links by P/N (23 Oct, \eref{PDM-04}) & Work files are \file{<P/N>.FCStd} and links are rewritten to that name, so a link
  names its child item.\\
PCBA as one line (23 Oct, \eref{ECD-05}) & A part with \code{BomDetail} off is a leaf (\file{BomObject.cpp}, \code{[BW]}).\\
Approvals in FreeCAD for now (23 Oct, \eref{PDM-05}) & Ops takes over approvals later; today Ops only needs to
  \emph{receive} BOMs.\\
\bottomrule
\end{tabularx}
\end{center}

%----------------------------------------------------------------------
\subsubsection*{1. What the spike showed (24 Oct)}

A script (\file{pdm-plans/spike\_bom.py}, headless \code{FreeCADCmd} on a copy of your configuration) made five
files: a bolt, a plate and a riser as parts; a \emph{bracket} assembly with one plate and a link array of four bolts;
and a \emph{top} assembly with the riser, two brackets and two loose bolts. It imported them into a scratch folder vault
with the Create Vault engine (\file{vault\_setup.py}), revised the bracket twice, and compared three sources of truth:
the vault's model, the links in each checked-in file, and FreeCAD's Assembly BOM.
\begin{center}
\small
\begin{tabularx}{\linewidth}{@{}>{\hsize=1.15\hsize}X>{\hsize=0.85\hsize}X@{}}
\toprule
Check & Result\\
\midrule
Who calls \code{set\_bom} in the product code (\code{grep}) & nobody: every vault's BOM is empty today\\
\code{release\_check.where\_used} (new with \eref{PDM-05}'s Retire) & item-level: latest file's links plus BOM
  lines, through sub-assemblies; one revision (the latest) per user\\
Vault model after \code{set\_bom} of the top and the bracket & explode: bolt 10, plate 2, bracket 2, riser 1 (correct)\\
Bracket's BOM replaced by one without bolts & where-used of the bolt drops the bracket \emph{and} the top: the model keeps
  one BOM per item, so the earlier revision's use is lost\\
Revision of the bracket after two check-ins & still \code{01} (three versions); a check-in makes a version, not a revision\\
Links read from each stored version (\code{linked\_files}) & child P/Ns per version; 7 versions read in 0.6\,ms\\
Bracket with the bolt array deleted headless & 4 orphan \code{App::LinkElement}s stay and still carry an XLink to the
  bolt: the raw links \emph{over-report}\\
Assembly BOM, headless (\code{Assembly::BomObject} in a BOM group) & works in \code{FreeCADCmd}\\
\quad top level only & riser 1, bracket 2, bolt 2 (correct)\\
\quad indented & bracket 2, plate 1, \textbf{bolt 6} under the bracket, \textbf{no top-level bolt row}\\
Ops (\file{bruceworks}, local \code{main}) & one route, \code{/api/health}; the dashboard's CAD-register tile is sample
  data; no BOM endpoint\\
\bottomrule
\end{tabularx}
\end{center}

\paragraph{Three findings to act on.}
\begin{enumerate}[nosep]
  \item \textbf{BOMs must be kept per revision.} \code{Vault.set\_bom(parent, lines)} replaces the item's BOM, so
        where-used can only describe today. The acceptance (``every assembly revision'') needs the BOM of each revision
        kept, and the released one frozen with its file, as \code{released\_files} already does for the file.
  \item \textbf{The BOM comes from the assembly, not from raw links.} Links are right for fetching the files an assembly
        needs (\eref{PDM-04}), but they include stale references and say nothing about quantities. The same applies to
        the release checklist's component line (\file{release\_check.py}, \code{components\_of}) and to
        \code{release\_check.where\_used}, which Retire\ldots{} uses: both read the \emph{latest} file's links, so they
        over-report stale links and cannot name an older revision.
  \item \textbf{The indented Assembly BOM merges across levels.} In \code{BomObject::addObjectChildrenToBom} the loop that
        adds a repeated sibling to its first row scans every row since the first sibling, which includes the rows of an
        earlier sibling's children. The top's two loose bolts were added to the bolt row inside the bracket (4 + 2 = 6)
        and their own row vanished. The fix keeps a list of this level's rows and searches only that (about ten lines,
        \code{[BW]}, and an upstream candidate: upstream main has the same loop).
\end{enumerate}

%----------------------------------------------------------------------
\subsubsection*{2. Theory: the BOM graph over revisions}

\paragraph{The graph.} Let $R$ be the set of revisions, each written $p@r$ (item $p$ at revision $r$). A checked-in
assembly revision $a$ has BOM lines $(a, c, q_{ac}, n_{ac})$ to child items $c$ with quantity $q_{ac}$ and find number
$n_{ac}$. A line names an \emph{item}, so to walk the graph each line must be \emph{resolved} to a child revision:
\[
  \rho(a, c) \;=\;
  \begin{cases}
    \text{the revision of } c \text{ recorded when } a \text{ was checked in or released} & \text{(resolved, ``as built'')}\\
    \text{the latest released revision of } c \text{ now} & \text{(floating, ``as current'')}
  \end{cases}
\]
With $\rho$ fixed, the BOM is a directed graph $G = (R, E)$ with an edge $a \to \rho(a,c)$ weighted $q_{ac}$. FreeCAD
allows two files to link each other, but a BOM cannot contain itself, so $G$ must be acyclic; the reference model
already refuses a cycle (\code{\_has\_cycle}, a depth-first search with three colours) and the check-in will report the
cycle's path.

\paragraph{Roll-up.} The total quantity of revision $v$ under the top revision $t$ is the sum over paths of the product
of quantities (as in \S\ref{sec:asm}),
\[
  Q_t(v) = \sum_{\pi:\,t \rightsquigarrow v}\ \prod_{(x,y)\in\pi} q_{xy},
  \qquad
  Q(t)=1,\quad Q(y) = \sum_{x \to y} Q(x)\,q_{xy}\ \text{in topological order,}
\]
which costs $O(|R|+|E|)$. The \emph{indented} BOM lists the paths with the per-parent quantity $q$ on each row; the
\emph{flattened} BOM lists $Q_t(v)$ for the leaves. In the spike, $Q(\text{bolt}) = 2\cdot4 + 2 = 10$. Leaves are parts,
COTS items and any part with \code{BomDetail} off. Quantities are counts for now (question~6).

\paragraph{Where-used as reverse reachability.} Where-used of revision $v$ is the set of revisions from which $v$ can be
reached,
\[
  W(v) = \{\,a \in R : a \rightsquigarrow v\,\},
  \qquad
  W(p) = \bigcup_{r} W(p@r)\ \text{for an item } p,
\]
found by a breadth-first walk over the reversed edges with a visited set, $O(|R|+|E|)$. For a few thousand revisions
this is milliseconds, so it is computed from the journal on demand, with no index to keep in step. Each answer carries
the path and $Q_a(v)$, so ``\code{LE-02001} A uses 10'' can be expanded to ``via the bracket, $2\times4$, and 2 direct''.

\begin{figure}[!htbp]
\centering
\begin{tikzpicture}[>=Stealth,x=2.9cm,y=-1.6cm,
  r/.style={draw=black!35,fill=white,rounded corners=2pt,font=\scriptsize,minimum width=2.2cm,minimum height=0.6cm,align=center},
  hit/.style={r,draw=bwred,fill=bwred!8,line width=0.8pt},
  q/.style={font=\tiny,fill=white,inner sep=1pt},
  e/.style={->,black!45},
  he/.style={->,bwred,line width=0.9pt}]
\node[hit] (pa) at (0.55,0) {Pedestal \textbf{A}\\[-1pt]\tiny released};
\node[r]   (pb) at (2.0,0) {Pedestal \textbf{B}\\[-1pt]\tiny released (ECO)};
\node[hit] (ka) at (0,1) {Bracket \textbf{A}};
\node[r]   (ri) at (1.25,1) {Riser A};
\node[r]   (kb) at (2.5,1) {Bracket \textbf{B}};
\node[hit] (bo) at (0,2) {\textbf{Bolt M5$\times$12} A\\[-1pt]\tiny COTS};
\node[r]   (pl) at (1.25,2) {Plate A};
\node[r]   (rv) at (2.5,2) {Rivet A\\[-1pt]\tiny COTS};
\draw[he] (pa) -- node[q]{$\times2$} (ka);
\draw[he] (pa.west) -- ++(-1.25,0) |- node[q,pos=0.25]{$\times2$} (bo.west);
\draw[e] (pa) -- node[q]{$\times1$} (ri);
\draw[e] (pb) -- node[q]{$\times1$} (ri);
\draw[e] (pb) -- node[q]{$\times2$} (kb);
\draw[he] (ka) -- node[q]{$\times4$} (bo);
\draw[e] (ka) -- node[q]{$\times1$} (pl);
\draw[e] (kb) -- node[q]{$\times1$} (pl);
\draw[e] (kb) -- node[q]{$\times4$} (rv);
\end{tikzpicture}

\vspace{6pt}
\footnotesize
\begin{tabularx}{\linewidth}{@{}>{\raggedright\arraybackslash}p{5.2cm}Xl@{}}
\toprule
How the graph is kept & Where-used of the bolt & Verdict\\
\midrule
One BOM per item, latest wins (today's model) & nothing (the bracket's BOM is now B, with rivets) & misses Pedestal A\\
Per revision, lines to items, not resolved & Bracket A $\rightarrow$ ``Bracket'' $\rightarrow$ Pedestal A \emph{and B} & reports B wrongly\\
Per revision, lines resolved (proposed) & Bracket A (4), Pedestal A (10) & exact\\
\bottomrule
\end{tabularx}
\caption{Where-used as reverse reachability over revisions (example). An ECO replaced the bracket's bolts with rivets:
Bracket~B and Pedestal~B no longer use the bolt. Walking the red edges backwards from the bolt finds exactly the
revisions that use it; quantities multiply along the path ($2\times4+2=10$).}
\label{fig:pdm06-whereused}
\end{figure}

\paragraph{Resolved or floating.} The data model (figure~\ref{fig:pdm-model}) lets a released assembly float to its
children's latest released revisions. Your numbering standard for Leonardo (LE-STD-0001) says the opposite in practice:
\emph{``if a child's released config changes an assembly's BOM, bump the parent rev''}, which is the resolved reading.
Both can be stored at no cost: each line records the child item \emph{and} the child revision seen at check-in, frozen at
release. Where-used then answers exactly (resolved), and can also show ``may pick up'' (floating) as a second list
(question~2).

%----------------------------------------------------------------------
\subsubsection*{3. Design}

\begin{figure}[!htbp]
\centering
\begin{tikzpicture}[>=Stealth,node distance=4mm and 6mm,
  b/.style={draw=bwblue!70,fill=bwblue!5,rounded corners=2pt,font=\scriptsize,align=center,minimum height=0.95cm,text width=2.05cm},
  g/.style={b,draw=black!30,fill=black!3},
  t/.style={font=\tiny\color{bwgrey},align=center}]
\node[b] (ci) {Check in an assembly (\eref{PDM-04})};
\node[b,right=of ci] (bom) {BOM walker: parts only, per-parent qty};
\node[b,right=of bom] (j) {Journal: \code{set\_bom} with revision, version, resolved children};
\node[b,right=of j] (rel) {Release freezes the version's BOM (\eref{PDM-05})};
\node[g,right=of rel] (ops) {Outbox $\rightarrow$ Ops (when its API exists)};
\draw[->] (ci) -- (bom); \draw[->] (bom) -- (j); \draw[->] (j) -- (rel); \draw[->] (rel) -- (ops);
\node[t,below=1mm of bom] {refuses a cycle;\\names children not in the vault};
\node[t,below=1mm of j] {older lines (child, qty)\\still replay};
\node[t,below=1mm of ops] {idempotent by\\(P/N, rev, SHA-256)};
\end{tikzpicture}
\caption{Where a revision's BOM is made, kept and sent.}
\label{fig:pdm06-flow}
\end{figure}

\begin{itemize}[nosep]
  \item \textbf{The BOM walker} (\file{bruceworx\_bom.py}, FreeCAD side) walks an open assembly the way
        \code{BomObject} does (links, link arrays counting visible elements, \code{AssemblyLink}, \code{BomDetail}) but
        returns structured rows: P/N (the child's \code{PartNumber}, else its file name), quantity per parent, find
        number, and the child's revision from the vault. It keeps only items: a part's bodies are not rows. It is the
        same walker \eref{ASM-04} needs for its indented, flattened and CSV views and \eref{DRW-02} for the parts list,
        so it is built once, under ASM-04 if that comes first. The \code{BomObject} merge bug is fixed alongside, because
        that spreadsheet is what you see in the tree.
  \item \textbf{Check-in.} When the file is an assembly, the walker runs on the open document before the file is
        stored. A child that is not a vault item is listed with \emph{Import\ldots} (the check-in panel already offers it
        for links), and a cycle stops the check-in with its path. The lines go into the journal with the version they
        describe.
  \item \textbf{Journal and model.} \code{set\_bom} gains \code{revision}, \code{sha256} and per line the child revision
        and find number; lines with only (child, qty) from older journals still replay. \code{pdm.Vault} keeps
        \code{bom[(parent, revision)]}, with the BOM of the version that was released frozen beside
        \code{released\_files}. \code{explode}, \code{flattened} and \code{where\_used} take a revision, and
        \code{where\_used} returns (P/N, revision, total quantity, path). \code{release\_check.where\_used} is not replaced
        but widened: it keeps its rows and its \code{via} column (``BOM'', ``file link''), gains the revision, the
        quantity and the path, reads the per-revision BOM first, and falls back to file links only for versions that have
        no BOM yet (marked ``file link only''). Retire\ldots{} keeps using it, so its rule (no retiring a part a released
        assembly still uses) becomes per revision with no change to the panel. The OpenAPI contract changes the same way:
        \code{/items/\{pn\}/revisions/\{rev\}/bom} and a where-used that returns revisions.
  \item \textbf{Release.} The release checklist's component line reads the BOM lines instead of raw links (finding~2).
        Releasing an assembly freezes its BOM; an ECO that releases a child does not change a parent's released BOM.
  \item \textbf{In the window.} Right-click an item in the Vault panel or the Browser: \emph{Where used\ldots} opens a
        list of assembly revisions with state chips, the path and the quantity (toggle: \emph{exact} or \emph{may pick
        up}); \emph{Bill of materials\ldots} shows a revision's BOM as indented or flattened, with \emph{Export CSV}.
  \item \textbf{Ops.} Each release appends one line per BOM to \file{outbox/<user>@<mac>.jsonl} in the vault (one
        writer per file, the journal's rule). When the Ops API exists (\eref{PDM-03}'s reader, bearer token from
        \code{BW\_OPS\_TOKEN}), a sender posts the lines it has not posted yet; Ops stores them by (P/N, revision,
        SHA-256), so sending twice changes nothing. Until then the outbox is the record, and \emph{Export CSV} covers
        anyone who needs the BOM now.
\end{itemize}

%----------------------------------------------------------------------
\subsubsection*{4. Build order and effort}

\begin{center}
\small
\begin{tabularx}{\linewidth}{@{}lXl@{}}
\toprule
\# & Step & Effort\\
\midrule
1 & Reference model: BOM per revision, resolved children, explode, flattened and where-used per revision; contract tests
  (FreeCAD-free) & 1\,d\\
2 & Journal format with old lines still replaying; folder tests; OpenAPI update & 0.5\,d\\
3 & BOM walker (shared with \eref{ASM-04}) and the \code{[BW]} fix of the indented merge, with a test on the spike's
  files & 1.5\,d\\
4 & Check-in hook: lines from the open assembly, children not in the vault, cycle report & 1\,d\\
5 & Release: freeze the BOM, component line from BOM lines & 0.5\,d\\
6 & Widen \code{release\_check.where\_used}; Where used\ldots{} and Bill of materials\ldots{} panels, CSV export & 1.5\,d\\
7 & Ops outbox and sender (off until Ops answers) & 1\,d\\
8 & Acceptance: the bolt across Pedestal A and B (figure~\ref{fig:pdm06-whereused}), then on \code{LE-02001};
  write-up & 0.5\,d\\
\midrule
 & Total & 7.5\,d\\
\bottomrule
\end{tabularx}
\end{center}
Steps 1--2 need nothing else and can start at once; step~3 is shared with \eref{ASM-04}, so whichever epic comes first
carries it.

\paragraph{Risks.} \emph{Re-reading old versions}: versions checked in before this epic have no BOM; a one-off
\emph{Rebuild BOMs} runs the walker in \code{FreeCADCmd} over each stored assembly version (opening every linked part, so
slow on large trees, but once). \emph{Configurations} (\eref{MOD-03}) will add a configuration to each line; the line
format leaves room for it. \emph{Ops} has no API yet, so the push cannot be tested end to end in this epic; the outbox
keeps that from blocking it. The \code{BomObject} fix touches an upstream file (one \code{[BW]} block).

%----------------------------------------------------------------------
\paragraph{Your answers to the questions below (24 Oct).} All seven as recommended. 1: where-used is \textbf{exact} by
default, with \emph{may pick up} as a second list in the same panel. 2: children are \textbf{pinned at release} (the
child revision is recorded) and the parent shows ``newer child released'', as LE-STD-0001 has it. 3: Ops receives
\textbf{released BOMs only} (first release and each ECO release). 4: the \textbf{outbox and Export CSV}; no Ops change
in this epic. 5: \textbf{every BOM line is a vault item}, bought parts included; check-in lists a child that is not an
item and release refuses one. 6: \textbf{counts} now, units later. 7: the indented merge in \file{BomObject.cpp} is
fixed now (\code{[BW]}) and offered upstream, the spike's files as the test. Building started 24 Oct.

\subsubsection*{Questions before PDM-06 starts}

\begin{enumerate}
  \item \textbf{What where-used answers.} Exactly the assembly revisions that were built with the part (resolved), or
        also those that could pick it up because they float to the latest released child?\\
        \emph{Recommendation:} exact by default, with \emph{may pick up} as a second list in the same panel.
  \item \textbf{Released assemblies: pinned or floating children.} The data model floats to the latest released child;
        LE-STD-0001 bumps the parent when a child change alters its BOM, which is pinning.\\
        \emph{Recommendation:} pin at release (record the child revision), and show ``newer child released'' on the
        parent so you decide whether it needs an ECO. This matches LE-STD-0001.
  \item \textbf{What Ops receives, and when.} Every check-in's BOM, or released BOMs only?\\
        \emph{Recommendation:} released BOMs only (first release and each ECO release); WIP BOMs stay in the vault.
  \item \textbf{Until Ops has an API.} Is the outbox in the vault plus \emph{Export CSV} enough, or do you want a file
        dropped somewhere Ops already reads?\\
        \emph{Recommendation:} outbox and CSV; no Ops change in this epic.
  \item \textbf{Hardware as items.} The acceptance needs the bolt to be a vault item. Should every BOM line be a vault
        item, including bought parts (a COTS item, \code{09xxx} in LE-STD-0001)?\\
        \emph{Recommendation:} yes; check-in lists any child that is not an item, and release refuses one.
  \item \textbf{Quantities.} Counts only, or units too (cable by length, adhesive by mass)?\\
        \emph{Recommendation:} counts now; a unit and quantity-per property later, when the first harness needs it.
  \item \textbf{The Assembly BOM bug.} Fix the indented merge in \file{BomObject.cpp} now (\code{[BW]}) and offer it
        upstream?\\
        \emph{Recommendation:} yes, with the spike's files as the test.
\end{enumerate}
